\def\ChapterCode{SAN}
\chapter
[\CO{[SAN]} Sanitätstechniken]
{Sanitätstechniken}
% \AasRsN{RS~XIII}
\label{chp:SAN}

\ChapterIntro
{%
Dieses Kapitel behandelt diverse allgemeine und
spezielle Maßnahmen der Sanitätshilfe. 
}
{%
\begin{description}
  \item[Maintainer:] Diverse
  \item[Autoren:] Diverse
  \item[Reviewer:] Standard-Reviewprozess
  \item[Version:] {\VarDocVersionTypeName} ({\VarDocVersionTypeDescription})
  \item[SHA1:] {\DocGitCommit}
\end{description}

{\TxTeilkapitelAASS}
}

% \author{Gabriel, Hochreiter, Motal}

% \minitoc
% \includegraphics[angle=270,width=\linewidth]{./images/leitsystem/Leitschema-PAM_SAN.pdf}

\clearpage % TEMP temp clearpage
% \TextSlide
% {\TxQuerverweise}
% {~\par\noindent\includegraphics[angle=270,width=\linewidth]{./images/leitsystem/Leitschema-PAM_SAN.pdf}}
% {\begin{itemize}
%   \item \EE{Medizinprodukte allgemein}: \CO{[MP1]} (\ref{chp:MP1})
%   \item \EE{Untersuchungen}: \CO{[BAU]} (\ref{chp:BAU})
% %   \item \EE{Sanitätshilfemaßnahmen}:  \CO{[SAN]} (\ref{chp:SAN})
% %   \item[Anamnese und Patientengespräch:]  Kapitel [AUP] (\ref{chp:AUP})
% \end{itemize}}
% {}
% {}
% {}
% % 
% % \Text{Inhalt}{%
% % Dieses Kapitel behandelt die Durchführung von Maßnahmen der Sanitätshilfe.
% % }{}{}{}{}
% % 
% % 
% % \Text{{\TxQuerverweise}}{%
% % \begin{compactdesc}
% %   \item[Beschreibung der Medizinprodukte:] Kapitel [MP1, MP2] 
% %   (\ref{chp:MP1}, \ref{chp:MP2}).
% %   \item[Untersuchungen:] Kapitel [BAU] (\ref{chp:BAU}).
% %   \item[Sanitätshilfemaßnahmen:]  Kapitel [SAN] (\ref{chp:SAN}).
% % %   \item[Anamnese und Patientengespräch:]  Kapitel [AUP] (\ref{chp:AUP}).
% % \end{compactdesc}
% % }{}{}{}{}


% \clearpage
\section{Die \HeadingKeyword{Hygienische Händedesinfektion} soll die Infektionskette unterbrechen}
\label{topic:hygienische-haendedesinfektion}

\Text
{\TxBeschreibung}
{Als hygienische Händedesinfektion wird die Desinfektion der
Hände und Handgelenke mittels einer Desinfektionslösung und
standardisierter Handbewegungen bezeichnet. Sie ist eine
Routinemaßnahme.
\EE{Ziel} ist  die 
\EE{Keimreduktion} und damit die \EE{Vermeidung von
Kreuzinfektionen} (\EEE{{Unterbrechung der
        Infektionskette}}). Die normale Händewaschung
        entfernt nur groben Schmutz und sonstige mechanische Verunreinigungen. Die
Desinfektionslösung reduziert jedoch viele
Keime aktiv und gezielt.

Übungen mit Fingerfarben
zeigen, dass mit einer herkömmlichen Handwaschtechnik große
Teile der Hand nicht mit Desinfektionsmittel
benetzt werden (Fingerzwischenräume,
Handgelenk, Daumenkuppe, \ldots). Abhilfe schaffen die im
Folgenden beschriebenen standardisierten Handbewegungen.}
{}
{}
{}
{}

\begin{Synopsis}
  \SynopsisItem{\I{Kreuzinfektion}: Patient~1 $\rightarrow$ Personal $\rightarrow$ Patient~2}
  \SynopsisItem{Die Hände des Personals sind
          die häufigsten Keimüberträger.}
  \SynopsisItem{Die hygienische Händedesinfektion ist \EE{vor} bzw. \EE{nach
  jedem Patientenkontakt} durchzuführen!}
\end{Synopsis}



\begin{PictureSeriesWide}{}{\EE{Wozu Händehygiene?} Multiresistente
Keime und ihre Auswirkungen in der Praxis. 
\SourcePicture{Ch. Gabriel/KARPAT}}
\PictureSeriesRowWide{3}
{./images/gabriel-christine-aass/staph-aureus.jpg}
{Auf einem Nährboden wurde eine Bakterienkultur angelegt.
Jedes der sechs runden Plättchen enthält ein Antibiotikum.
Man kann sehen dass in der Umgebung der Plättchen keine
Bakterien
wachsen.}
{./images/gabriel-christine-aass/staph-aureus-mrsa.jpg}
{Diese Bakterien (gleicher Keim!) ignorieren 4 von 6
Antibiotika-Plättchen, d.\,h. die entsprechenden Antibiotika wirken
nicht gegen diesen Keim. Er ist \EE{multiresistent}. Die
beste Maßmnahme dagegen: \Ca{Gar nicht erst übertragen!}}
{./images/gabriel-christine-aass/amputat-us1.jpg}
{Oft die Folge: Amputation eines Unterschenkels mit einer nicht
abheilenden, mit einem muliresistenten
Staph. aureus (MRSA) infizierten Wunde.}
{empty}
{}
\end{PictureSeriesWide}


\Text
{Material}
{\begin{compactitem}
    \item Alkoholische Händedesinfektionslösung, z.\,B.:
    Sterilium\texttrademark{}, Manopronto\texttrademark{}
    o.\,ä.
\end{compactitem}}
{}
{}
{}
{}



% \end{multicols}
% \end{WidePar}


\TextSlide
{Handlungsschritte}
{Siehe \FullPrettyRef{fig:haendedesinfektion}.


\begin{PictureSeriesWide}{}{Hygienische Händedesinfektion. \EE{Einwirkzeit beachten!} (lt.
Herstellerangaben, i.\,d.\,R. mind. 30 Sekunden)
\SourcePicture{Blacky}\label{fig:haendedesinfektion}}
\PictureSeriesRowWide{3}
{./images/hygiene/haendedesinfektion01.png}
{Eine Portion alkoholisches
    Händedesinfektionsmittel (ca. 3\Ml* = 1~Hohlhand)
    mittels Ellenbogen aus dem Spen\-der entnehmen; Hände
    und Handgelenke vollständig benetzen}
{./images/hygiene/haendedesinfektion02.png}
{Handinnenflächen gegeneinander reiben, geschlossene Finger.}
{./images/hygiene/haendedesinfektion03.png}
{Handgelenke einreiben}
{empty}
{Mit ineinander verschränkten Fingern Handinnenflächen
    gegeneinander reiben}
%
\PictureSeriesRowWide{3}
{./images/hygiene/haendedesinfektion04.png}
{Mit ineinander verschränkten Fingern Handinnenflächen
    gegeneinander reiben}
{./images/hygiene/haendedesinfektion05.png}
{Mit rechter Handinnenfläche linken Handrücken reiben und
    umgekehrt, dabei Finger ineinander verschränken}
{./images/hygiene/haendedesinfektion06.png}
{Hände ineinander verhaken und Finger gegeneinander bewegen}
{empty}
{Daumen mit gegenüberliegender Hand vollständig umschließen
    und rotierend reiben. Daumenkuppen nicht vergessen!} 
%
\PictureSeriesRowWide{3}
{./images/hygiene/haendedesinfektion07.png}
{Daumen mit gegenüberliegender Hand vollständig umschließen
    und rotierend reiben. Daumenkuppen nicht vergessen!}
{./images/hygiene/haendedesinfektion08.png}
{Fingerkuppen im Handteller kreisförmig einreiben bis Alkohol
    verdunstet ist}
{empty}
{Raum für Phantasie}
{empty}
{}
\end{PictureSeriesWide}}
{\FullPrettyRef{fig:haendedesinfektion}}
{}
{}
{}




\EE{Einwirkzeit beachten!} (lt.
Herstellervorgaben, i.\,d.\,R. mind. 30 Sekunden)


\nocite{KirschnickPflege2}





\needspace{0.5\textheight}
\section{\HeadingKeyword{Hebe-} und \HeadingKeyword{Transporttechniken}}
\index{Hebetechniken}
\index{Transporttechniken}

\TODO{GABS}{2010-08-14}{Hebe- und Transporttechniken: Text
fehlt}{} 
% \Text{{\TxQuerverweise}}{%
% \begin{inparaitem}
%   \item Tragsessel: \prettyref{topic:tragsessel};
%   \item Fahrtrage, \enquote{Einheitskrankentrage}: \ref{topic:fahrtrage};
% \end{inparaitem}
% }{}{}{}{}

Das Fachpersonal hat während des gesamten Transports die
Verantwortung für den Patienten, dies gilt auch für gehfähige
Patienten. Sollte der Patient auf Grund seines hohen Alters
oder seiner Kreislaufsituation nicht selbstständig gehen
können, muss er entweder \EE{sitzend} oder \EE{liegend}
transportiert werden. 
Patienten, bei denen mit einer Zustandsverschlechterung zu
rechnen ist, sollen grundsätzlich liegend
transportiert werden, da ein Umlagern im Fahrzeug sehr
umständlich und zeitraubend ist.


\Experimental{
\subsubsection{\HeadingKeyword{Umlagern} von Patienten}
\index{Umlagern}
\label{topic:umlagern}

{\TakeNotesLarge}}



\Experimental{
\subsubsection{Umgang mit einem \HeadingKeyword{Rollstuhl}}
\index{Rollstuhl}
\label{topic:rollstuhl}

% \prettyref{fig:rollstuhl-umgang}

{\TakeNotesLarge}

}





\subsection{\HeadingKeyword{Tragering}}

\label{topic:tragering}

Das Tragen mit dem Tragering funktioniert nur bei bewusstseinsklaren und
kooperativen Patienten, außerdem sollte das Stiegenhaus, oder die zurück
zulegende Strecke, breit genug sein. Die Technik eignet sich nur für kurze
Distanzen.

\begin{PictureSeriesWide}{}{Bilderserie: Transport mit einem
\EE{Tragering} \SourcePicture{Motal}}
\PictureSeriesRowWide{3}
{./images/motal-aass/00800/tragering1b.jpg}
{Tragering}
{./images/motal-aass/00800/tragering2b.jpg}
{Transport mit dem Tragering}
{./images/motal-aass/00800/tragering3b.jpg}
{Transport mit dem Tragering}
{}
{}
\end{PictureSeriesWide}






\subsection{\HeadingKeyword{Tragetuch}}

\index{Tragetuch}
\label{topic:bergetuch}
\label{topic:tragetuch-taetigkeit}

Das Tragetuch (ESE-Tuch) eignet sich ideal für den Transport
von liegenden Patienten bei denen der Einsatz einer
Krankentrage eventuell aus Platzgründen nicht möglich ist. Es wird auch
oft zum Umlagern von Patienten oder zum Retten aus einer
Gefahrenzone verwendet. 
Der Transport mit Hilfe des Tragetuchs ist \EE{nicht
wirbelsäulenschonend}. Ist eine Immobilisation der
Wirbelsäule notwendig, soll stattdessen eine andere Technik
(Vakuummatratze, etc.) angewendet werden.

Die Füße des Patienten zeigen
talwärts, in der Ebene wird der Patient wird immer mit den
\EE{Füßen voran} getragen.



\begin{PictureSeriesWide}{}{Bilderserie: Transport mit dem Tragetuch\index{Tragetuch!Bilderserie}}
\PictureSeriesRowWide{3}
{./images/motal-aass/00800/tragetuch1b.jpg}%1
{Tuch aufbreiten \ldots}
{./images/motal-aass/00800/tragetuch2b.jpg}%2
{{\ldots} die näher zum Patient liegende Hälfte
umschlagen\ldots }
{./images/motal-aass/00800/tragetuch3b.jpg}%3
{{\ldots} die obere Hälfte nochmals
umschlagen\ldots }
{empty}%4
{}
%
\PictureSeriesRowWide{3}
{./images/motal-aass/00800/tragetuch4b.jpg}%1
{{\ldots} und das Tuch an den Patienten legen.}
{./images/motal-aass/00800/tragetuch5b.jpg}%2
{Patient leicht anheben und Tuch bis zur Hälfte
unterschieben \ldots}
{./images/motal-aass/00800/tragetuch6b.jpg}%3
{{\ldots} Patient schonend zurück legen.}
{empty}%4
{}
%
\PictureSeriesRowWide{3}
{./images/motal-aass/00800/tragetuch7b.jpg}%1
{Von der anderen Seite ebenfalls leicht anheben und
Tuch ausbreiten. }
{./images/motal-aass/00800/tragetuch8b.jpg}%2
{Zu dritt das Tuch fest anpacken \ldots }
{./images/motal-aass/00800/tragetuch9b.jpg}%3
{\ldots und auf Kommando gemeinsam anheben}
{empty}%4
{}
\end{PictureSeriesWide}



\clearpage
\subsection{\HeadingKeyword{Tragsessel}}



\TextSlidePicture
{Umgang mit dem Tragsessel}%1 Titel
{\label{topic:tragsessel}%
Der Tragsessel ist eines der am häufigsten verwendeten
Transportmittel für Patienten im Krankentransport. In der
Notfallrettung wird er häufig verwendet, um
bewusstseinsklare Patienten zum Fahrzeug zu
transportieren. Es gibt viele verschiedene Modelle, welche
unterschiedliche Funktionen aufweisen können.

\E{Notfallpatienten} dürfen \E{nur im begründeten
Ausnahmefall} während der Fahrt im Tragsessel belassen
werden, da bei einer Verschlechterung des Zustandes das
Umlagern des Patienten im Fahrzeug sehr kompliziert (bei
bewusstlosen Patienten nahezu unmöglich) ist. 

Der Patient ist im Tragsessel \EE{immer anzuschnallen}. Beim
Transport über Stiegen hat der Patient immer \EE{talwärts} zu blicken.}%2 Text
{}%3 Slide
{./images/pallinger-christoph-aass/Tragsessel_32966-prelim-edited2-crop2-AASS-0112mm}%4 Bilddatei
{Tragsessel}%5 Bildtitel
{}%6 Bildbeschreibung
{Ch. Pallinger}%7 Quelle
{MfG}%8 Lizenz
% Vormals \Layout{20}

\Experimental{\Text{Zutragestuhl}
{\label{topic:zutragestuhl-mp}%
{\TakeNotesLarge}
}{

}
{}
{}
{}}



\subsection{Umgang mit \HeadingKeyword{Krankentragen}}
\label{topic:fahrtrage}

\Text{\TxBeschreibung}
{%
\index{Trage}%
\index{Krankentrage}%
\index{Fahrtrage}%
\index{Stollenwerk{\texttrademark}}%
\index{Ferno!Fahrtrage}%
%
Eine Krankentrage besteht aus einer Liegefläche für den
Patienten und 4 Holmen, mit welchen die Trage von zwei oder
mehr Personen getragen werden kann. Moderne Tragen verfügen
zudem über ein höhenverstellbares, abtrennbares Fahrgestell. 
Tragen verfügen über unterschiedliche \EE{Gurtensyteme},
grundsätzlich ist der Patient immer gemäß der
Herstellerangaben anzuschnallen (Schultern, Becken und
Beine!). Der \EE{Beingurt} ist sehr wichtig: Er verhindert
bei einem Frontalzusammenstoß, dass sich der Unterkörper des
Patienten überschlägt und die Wirbelsäule massive
Schädigungen erleidet.

Beim Transport über Stiegen hat der Patient immer
\E{talwärts} zu blicken.
Wenn ein Patient auf der Trage liegt, darf immer nur
\E{\EE{ein} Ende} gleichzeitig gehoben oder gesenkt werden,
da sonst die Gefahr des seitlichen Kippens besteht.
}{

}
{}
{}
{}

\begin{Synopsis}
  \SynopsisItem{Bei Transport in Fahrzeugen verhindert der
  \EE{Beingurt} schwere Verletzungen bei
  Frontalzusammenstößen!}
\end{Synopsis}
\begin{Danger}
  \item Der Beingurt ist bei Beförderung in Fahrzeugen unbedingt anzulegen!
\end{Danger}
% \subsubsection{Krankentrage mit Fahrgestell der Fa. Stollenwerk}
% 
% 

% \Text{\TxBeschreibung}
% {%
% \begin{compactdesc}
%     \item[Hersteller:] Stollenwerk
%     \item[Produktfamilie:] Krankentrage, Fahrgestell
%     \item[Modell/Typ:] Trage 3002, 3003, 3006; Fahrgestell 2870
%     \item[Modellbesonderheiten:] \begin{inparadesc}
%     \item[Trage 3002:] Kopfhochstellung. 
%     \item[Trage 3003:] Kopf- und Fußhochstellung. 
%     \item[Trage 3006:] Kopf- und Fußhochstellung, Bauchdecken-Entlastung. 
%     \item[Alle Tragen:] sind mit allen Stollenwerk-Fahrgestellen kompatibel!
% \end{inparadesc}   
%     \item[Max. Belastbarkeit:] ---\Kg*
%     \item[Homepage:] \url{http://www.stollenwerk-koeln.de}
% %     \item[Hygienische Wiederaufbereitung:]
% \end{compactdesc}
% % {\TakeNotesLarge}
% }{}{}{}{}


\PictureWide
{Krankentrage der Firma Stollenwerk{\texttrademark} mit
Fahrgestell und einer Auflage der Firma
Schnitzler{\texttrademark} mit Brust-, Bein- und
Rucksackgurten}%1
{}%2
{}%3
{Ch. Pallinger}%4
{MfG}%5
{./images/pallinger-christoph-aass/Trage_32970_v2-AASS-0176mm.jpg}%6
{}%7
{}%8






% \subsubsection{Krankentrage mit Fahrgestell der Fa. Ferno}


\clearpage % TEMP temp clearpage
\section{\HeadingKeyword{Lagerung}}
\subsection{\HeadingKeyword{Lagerungsarten}}
\label{topic:lagerungsarten}
\index{Lagerungsarten}


\begin{PictureSeriesWide}{}{Bilderserie: Lagerungsarten}
\PictureSeriesRowWide{3}
{./images/motal-aass/00800/lagerungflachb.jpg}
{Flachlagerung}
{./images/motal-aass/00800/lagerung30gradb.jpg}
{Lagerung mit 30\Deg erhöhtem Oberkörper}
{./images/motal-aass/00800/lagerungoberkoerperhochb.jpg}
{Lagerung mit stark erhöhtem Oberkörper }
{empty}
{}
%
\PictureSeriesRowWide{3}
{./images/motal-aass/00800/lagerungseiteb.jpg}
{Seitenlagerung}
{./images/motal-aass/00800/lagerungkollaps2b.jpg}
{Lagerung mit erhöhten Beinen.}
{./images/motal-aass/00800/lagerungkollaps1b.jpg}
{Lagerung mit
erhöhten Beinen unter Schonung der Wirbelsäule und des
Beckens (bei Verdacht auf Wirbelsäulen- oder
Beckenverletzungen)}
{}
{}
%
\PictureSeriesRowWide{3}
{./images/motal-aass/00800/lagerungabdomenb.jpg}
{Lagerung zur Entspannung der Bauchdecke}
{./images/motal-aass/00800/lagerungthromboseb.jpg}
{Lagerung mit erhöhter Exxtremität}
{./images/motal-aass/00800/lagerungembolieb.jpg}
{Lagerung mit herabhängender Gliedmaße}
{}
{}
\end{PictureSeriesWide}

\Commentary[Lagerung mit erhöhten Beinen]{Diese Lagerung
wurde früher als in Hinblick auf den Hypovolämischen Schock als
\I{Schocklagerung} bezeichnet. Da sich die Lagerung je nach
Schockart unterscheidet (Die Lagerung mit erhöhten Beinen
ist beim kardiogenen Schock lebensgefährlich!), soll diese
Bezeichnung nicht
mehr verwendet werden.}


\clearpage




\needspace{0.5\textheight}
\section{Manuelle \HeadingKeyword{Traumatechniken}}
% Michael Motal

\Text{{\TxQuerverweise}}{%
\begin{itemize}
  \item Traumacheck: \prettyref{topic:traumacheck}, Seite
  \pageref{topic:traumacheck}.
\end{itemize}
}{}{}{}{}

\subsection{\HeadingKeyword{Helmabnahme} und \HeadingKeyword{manuelle HWS-Stabilisierung}}
\label{topic:helmabnahme}
\label{topic:hws-stabilisierung-manuell}
\index{Helmabnahme}
\index{HWS!manuelle Stabilisierung}


\Text{\TxBeschreibung}
{Ein Patient mit aufgesetztem Sturzhelm stellt eine 
besondere Herausforderung dar. Der Sturzhelm ist
grundsätzlich immer abzunehmen! }{}{}{}{}

\begin{ProcedureStepsNoHeading}
\Text{Material und Personal}{
\begin{compactitem}
\item 2 Helfer ({\color{HelferA}\E{hinterer} Helfer},
{\color{HelferB}\E{seitlicher} Helfer})
\item 1 HWS-Schiene 
\end{compactitem}
}{}{}{}{}


% \begin{multicols}{2}
\Text{Vorgehen}{
\begin{enumerate}
  \item \EE{Dem Patient
  in Blickrichtung nähern, um
  unnötige Bewegungen des Kopfes zu vermeiden, und
  gleichzeitig auffordern den Kopf nicht zu bewegen.}
    \item {\color{HelferA}\EE{Fixieren}: Der \E{hintere}
    Helfer fixiert den Helm von hinten und lässt nicht mehr
    los, bis der Helm entfernt ist. Zwischen Helm und
    Schritt eine Helmhöhe Abstand lassen, damit der Helm in
    einer Bewegung abgezogen werden kann.}
    \item {\color{HelferB}\EE{Ansprechen}: Der
    \E{seitliche} Helfer spricht den Patienten von vorne an, öffnet
    Visier und Kinnriemen.}
    \item {\color{HelferB}Der \E{seitliche} Helfer
    \EE{stabilisiert} den Kopf während der Helmabnahme: eine
    Hand fasst mit Daumen und Zeigefinger das Unterkiefer
    (\I{C-Griff}), die andere Hand stützt den Kopf von
    unten.}
    \item {\color{HelferA}\EE{Helm abziehen}: Der
    \E{hintere} Helfer zieht den Helm in einer
    \EE{S-förmigen Bewegung} ab: den Helm \EE{zuerst kippen
    bis die Nasenspitze} sichtbar ist, dann die Unterkante
    in der \EE{gegengleichen Bewegung} zum Helfer bewegen.}
    
    {\color{HelferB}Während der Helm abgezogen wird,
    rutscht der \E{seitliche} Helfer mit seiner Haltehand
    etwas nach.}
    \item {\color{HelferA}Der \E{hintere} Helfer zieht den
    Helm vorsichtig fertig ab und legt ihn sicher ab.}
    \item {\color{HelferA}\EE{Der \E{hintere} Helfer
    übernimmt den Kopf} im doppelten C-Griff: Daumen und
    Zeigefinger um ein Ohr. Der Kopf darf
    \EE{\E{nicht} auf Zug} gehalten werden, die \Abk{HWS}
    darf \EE{nicht gestaucht} werden.
    Der Kopf wird in
    Neutralposition
    stabilisiert. \cite{ITLSde6,AsboeLehrmeinungRd2011-07}
    \Commentary[HWS-Immobilisation / Zug]{ASBÖ: Nach der
    Lehrmeinung 07/2011 vom 01.02.2011
    \cite{AsboeLehrmeinungRd2011-07} darf der Kopf nicht mehr
    unter Zug gehalten werden. Dies ist konsistent mit
    \Abk{ITLS} (6.\,A.)
    \cite{ITLSde6}}}
    \item {\color{HelferB}Der \E{seitliche} Helfer legt
    eine \Abk{HWS}-Schiene an
    (\FullPrettyRef{topic:hws-schiene-anlegen}).}
\end{enumerate}
% \end{multicols}
}{}{}{}{}
\end{ProcedureStepsNoHeading}


\begin{PictureSeriesWide}{}{Bilderserie: Helmabnahme und Anlage
einer HWS-Schiene \SourcePicture{Motal}\label{fig:helmabnahme-hws-schiene}\label{fig:hws-schiene-anlegen}}
\PictureSeriesRowWide{3}
{./images/motal-aass/00800/helm1.jpg}%1
{Während des Annäherns den Patient auffordern, den Kopf nicht zu bewegen.}
{./images/motal-aass/00800/helm2.jpg}%2
{Den Kopf sofort mit den Händen fixieren, Visier öffnen und Patienten ansprechen.}
{./images/motal-aass/00800/helm3.jpg}%3
{Ggfs. Patient synchron und schonend auf den Rücken drehen.}
{empty}%4
{}
%
\PictureSeriesRowWide{3}
{./images/motal-aass/00800/helm4.jpg}%1
{Ggfs. Brille entfernen, Kinnriemen öffnen und behindernde
Kleidung entfernen.}
{./images/motal-aass/00800/helm5.jpg}%2
{Halswirbelsäule stabilisieren.}
{./images/motal-aass/00800/helm6.jpg}%3
{Helm kippen bis von oben die Nasenspitze sichtbar ist,
danach vorsichtig nach hinten wegziehen}
{empty}%4
{}
%
\PictureSeriesRowWide{3}
{./images/motal-aass/00800/helm7.jpg}%1
{Nach dem Abziehen Stabilisierung des Kopfes übernehmen. Keinen Zug und keine Stauchung ausüben!}
{./images/motal-aass/00800/helm8.jpg}%2
{Nach der Helmabnahme erfolgt sofort das Anlegen einer
HWS-Schiene. Der Kopf bleibt in Neutralposition zug- und
stauchungsfrei fixiert. Die passende Größe der HWS-Schiene
wird bestimmt und eingestellt.}
{./images/motal-aass/00800/hwsschiene3.jpg}%3
{HWS-Schiene zuerst unter dem Patienten durchführen, danach
am Kinn anlegen.}
{empty}%4
{}
\end{PictureSeriesWide}




\needspace{0.5\textheight}
\section{Ruhigstellung mittels \HeadingKeyword{Schienung}}
\label{topic:schienung}\label{topic:schienungsmaterial}

\subsection{\HeadingKeyword{HWS}-Immobilisationsschiene (Stifneck\texttrademark)}
\label{topic:hws-schiene-anlegen}\index{HWS-Schiene!anlegen}\index{Stifneck|see{HWS-Schiene}}
\label{topic:hws-schiene-beschreibung}
\index{HWS-Schiene!Beschreibung}

\HeadingSub{Gängige Marken: Stifneck{\texttrademark},
Stifneck select{\texttrademark},
Stifneck Paedi select{\texttrademark}}

% http://de.wikipedia.org/wiki/Cervicalstütze#cite_ref-1
% Gemäß Empfehlung der DGU | ref=[1] kann beim Fehlen folgender fünf Kriterien davon ausgegangen werden, dass keine instabile Wirbelsäulenverletzung vorliegt:
% Bewusstseinsstörung
% neurologisches Defizit
% Wirbelsäulenschmerzen oder Muskelhartspann
% Intoxikation
% Extremitätentrauma
% ↑ Langfassung der Leitlinie "Polytrauma / Schwerverletzten-Behandlung" Deutsche Gesellschaft für Unfallchirurgie, 2011, Seite 110


\Text{\TxBeschreibung}
{%
Die HWS-Schiene\ZFootnote{HWS-Schiene: Auch \I{HWS-Orthese}\index{Orthese!HWS}.} 
zum Ruhigstellen und Stabilisieren der {\bfseries
H}als{\bfseries W}irbel{\bfseries S}äule (HWS). Sie soll bei
jedem Verdacht auf eine Verletzung der (Hals-) Wirbelsäule
eingesetzt werden (\FullPrettyProcedureRef{YY14200B}); jedenfalls nach einer Helmabnahme.
Bei einem Schädel-Hirn-Trauma (\Abk{SHT}) kommt es oft zu Verletzungen 
der HWS, daher sollte hier die Indikation auch großzügig gestellt werden.

Es sind verschiedene
Fabrikate verfügbar (z.\,B. Stifneck\texttrademark, Ambu Perfit{\texttrademark}, \ldots). 
Es gibt
Schienen mit einer fixen Größe oder moderne Schienen, bei
welchen sich die gewünschte Größe einstellen lässt
(Stifneck select{\texttrademark} bzw. Stifneck Paedi
select{\texttrademark} für Kinder).

}{

}
{}
{}
{}

\PictureWide
{Einstellbare HWS-Schienen für Erwachsene und Kinder}%1
{Laerdal{\texttrademark} Stifneck Select{\texttrademark} und
Stifneck Pedi-Select{\texttrademark}}%2
{}%3
{Ch. Pallinger}%4
{MfG}%5
{./images/pallinger-christoph-aass/Stiffneck_32899_crop-sRGB-AASS-0112mm}%6
{}%7
{}%8


\needspace{0.5\textheight}
~
\begin{ProcedureStepsNoHeading}
\Text{Material und Personal}{
\begin{itemize}
    \item 2 Helfer ({\color{HelferA}\E{hinterer} Helfer},
{\color{HelferB}\E{seitlicher} Helfer})
    \item 1 \Abk{HWS}-Schiene (z.\,B. Stifneck\texttrademark)
\end{itemize}
}{}



% \begin{multicols}{2}
\Text{Vorgehen}{
Siehe auch \FullPrettyRef{fig:hws-schiene-anlegen}.
\begin{enumerate}
    \item \EE{Patient auffordern den Kopf nicht zu bewegen.}
    \item {\color{HelferA}Der hintere
    Helfer \EE{stabilisiert} den Kopf und die \Abk{HWS}
    in Neutralposition ohne Zug und ohne Stauchung.}
    \item {\color{HelferB}Beengende und behindernde
    \EE{Kleidung entfernen}.}
    \item {\color{HelferB}\EE{Größenmessung}: Der seitliche
    Helfer misst die Entfernung von Schulter bis Unterkiefer in Fingerbreiten
    aus.}
    \item {\color{HelferB}Stifneck entsprechend einstellen:
    gemessene Entfernung einstellen.
    
    Verschiebemechanismus des
    Stifneck blockieren, ab nun darf kein Verstellen mehr
    möglich  sein.}
    \item {\color{HelferB}Klettverschlussband so einfalten,
    dass es nicht am Boden schleifen kann (hält sonst nicht
    mehr).}
    \item {\color{HelferB}Nackenteil unter Patient durch
    führen (Vorsicht: Klettverschlussband).}
    \item {\color{HelferB}\Abk{HWS}-Schiene \EE{ am
    Kinn anlegen}.}
    \item {\color{HelferB}Fest schließen.}
    
Die Kinnspitze sollte ein kurzes Stück über Kante hinaus
schauen (kein Durchrutschen), allerdings nicht allzu weit
(falsch eingestellt).
\end{enumerate}
% \end{multicols}
}{}{}{}{}
\end{ProcedureStepsNoHeading}

\clearpage % TEMP temp clearpage
\subsection{\HeadingKeyword{Schaufeltrage}}
\label{topic:schaufeltrage-anwendung}
\label{topic:schaufeltrage-handhabung}
\index{Schaufeltrage}

\index{Schaufeltrage!Beschreibung}
\label{topic:schaufeltrage-beschreibung}

\HeadingSub{Gängige Marken: Ferno ScoopEXL{\texttrademark}}

\Text{\TxBeschreibung}
{%
Die Schaufeltrage wird primär für das \EE{Umlagern} von liegenden
Patienten mit dem Verdacht auf eine Verletzung der
\E{Wirbelsäule, Becken oder Oberschenkel} verwendet. Sie
ermöglicht eine möglichst bewegungsfreie und schmerzarme
Rettung bzw. Umlagerung. Das Anwendungsgebiet ist jedoch
sehr vielfältig. Es reicht vom sicheren Transport im
unwegsamen Gelände bis zum Retten aus einer LKW-Kabine. 

Letztendlich soll der Patient zwecks Schienung auf einer
\EE{Vakuummatratze} gelagert werden.
Je nach Fabrikat kann der Patient jedoch auch direkt auf
einer Schaufeltrage (ohne Verwendung einer Vakuummatratze)
immobilisiert werden, dazu werden spezielle Gurte und
Kopfstützen benötigt.

Die Schaufeltrage setzt sich aus 2 Längshälften zusammen;
die Fußteile sind längenverstellbar. Die Kopfteile haben
eine fixe Länge.
Sollte der Patient größer als die Trage sein, müssen Kopf
und Oberkörper auf jeden Fall trotzdem auf der Trage zu
liegen kommen. Ausnahme: Beinverletzung, bei der eine
Wirbelsäulenverletzung ausgeschlossen werden kann.

\Z{Das Spineboard
(\FullPrettyRef{topic:spineboard-beschreibung}) ist in
vielen Fällen eine Alternative zur Schaufeltrage
\cite{ITLSde6,PHTLS6,AsboeLehrmeinungRd2011-10}.}
}{

}
{}
{}
{}

\begin{figure}[H]
\subfloat[Eine Schaufeltrage Marke Ferno{\texttrademark} ScoopEXL{\texttrademark}]{\includegraphics[width=\linewidth]{./images/pallinger-christoph-aass/Taschen_32957_crop-AASS-0112mm}}

\subfloat[Schaufeltrage halb geöffnet]{\includegraphics[width=0.45\linewidth]{./images/pallinger-christoph-aass/Taschen_32958_crop-AASS-0112mm}}
\hfill
\subfloat[Schaufeltrage getrennt]{\includegraphics[width=0.45\linewidth]{./images/pallinger-christoph-aass/Taschen_32959_crop-AASS-0112mm}}

\Captionof{figure}{Schaufeltrage}
\end{figure}

Die Schaufeltrage kann vielfältig eingesetzt werden:

\begin{itemize}
  \item Retten und Aufschaufeln eines Patienten \dotfill
  \FullPrettyRef{topic:schaufeltrage-aufschaufeln}
  \item Umlagern eines Patienten \dotfill \FullPrettyRef{topic:schaufeltrage-umlagern}
  \item Immobilisation direkt auf der Schaufeltrage (manche
  Modelle) \dotfill \FullPrettyRef{topic:schaufeltrage-immobilisation}
\end{itemize}

\clearpage
\subsubsection{\HeadingKeyword{Aufschaufeln} mittels Schaufeltrage}
\label{topic:schaufeltrage-aufschaufeln}

\begin{ProcedureStepsNoHeading}
\Text{Material und Personal}{
\begin{itemize}
    \item Mindestens 2 Helfer, besser 3 Helfer
    \item Schaufeltrage
    \item Gurte für Schaufeltrage
\end{itemize}
}{}{}{}{}




% \begin{multicols}{2}
\Text{Vorgehen}{
\begin{enumerate}
    \item Schaufeltrage neben Patienten in zusammengebautem Zustand auf die
    richtige Länge einstellen. Sollte ein Patient zu groß sein,
    müssen Kopf und Oberkörper auf jeden Fall trotzdem auf der Trage
    liegen! 
    
    Vorsicht: Markierung der max. Länge bei Schraubmodell.
    \item \EE{Öffnen} der Trage: Schnappverschlüsse; je nach Methode nur den
    unteren oder beide öffnen.
    \item \EE{Aufschaufeln}:
    \begin{itemize}
      \item 1. Methode: Fußseite offen, Patient wird von oben
      kommend wie auf eine Schere aufgeladen, die Trage unter ihm
      geschlossen, dabei
      stabilisiert ein 3. Helfer
      den Kopf. Abschließend Kontrolle des sicheren Schlusses!
      
      Diese Methode ist besonders bei Becken- und
      Oberschenkelverletzungen sinnvoll und der zweiten
      Methode vorzuziehen da der Patient nicht am Becken
      bewegt wird.
      \item 2. Methode: Beide Teile sind getrennt.
      \begin{enumerate}
        \item Der Patient wird leicht an
        Schulter und Hüfte abwechselnd zur Seite gedreht, dabei
        stabilisiert ein 3. Helfer
        den Kopf und folgt der Drehbewegung achsengerecht.
        Die jeweiligen Tragenteile werden während der Drehbewegung
        unter den Patienten geschoben.
        \item Die Hälften werden zuerst am Kopf-, dann am Fußende
        geschlossen.
        \item Abschließend wird eine Kontrolle des sicheren Schlusses durchgeführt.
      \end{enumerate}
    \end{itemize}
    \item Angurten: mehrere Gurte\footnote{Anzahl der Gurte lt.
    Herstellerangabe. Wir empfehlen die Verwendung von mind. \E{3} Gurten};
    jeden Gurt zuerst um den Handgriff der einen, dann der anderen Hälfte
    führen,  dann über dem Patienten schließen und festziehen.
    \item Umlagern auf Vakuummatratze: Trage mit
    vorbereiteter Matratze
    (\prettyref{topic:vakuummatratze-anwendung}) neben
    Patienten stellen um diesen
    einen möglichst geringen Weg tragen/bewegen zu müssen.
    \item Gurte öffnen und entfernen.
    \item Hälften trennen: beide Schnappverschlüsse öffnen
    \item Hälften parallel zum Patienten flach wegziehen. Der Patient soll
    hierbei nicht bewegt werden!
    \item Sollten die Schnappschlösser klemmen, kann man die
    Trage wie bei der Scherentechnik an nur einer Seite
    öffnen und in V-Form (gleichzeitig, je ein Helfer auf
    einer Seite) entfernen.
\end{enumerate}
% \end{multicols}
}{}{}{}{}
\end{ProcedureStepsNoHeading}

\clearpage
\subsubsection{\HeadingKeyword{Umlagern} mittels Schaufeltrage}
\label{topic:schaufeltrage-umlagern}

\begin{ProcedureStepsNoHeading}
\Text{Material und Personal}{
\begin{itemize}
    \item Mind. 2 Helfer
    \item Patient auf Schaufeltrage
    (\FullPrettyRef{topic:schaufeltrage-aufschaufeln})
\end{itemize}
}{}{}{}{}




% \begin{multicols}{2}
\Text{Vorgehen}{
\begin{enumerate}
  \item Patient liegt, wie unter
  \FullPrettyRef{topic:schaufeltrage-aufschaufeln},
  beschrieben, auf der Schaufeltrage.
    \item Umlagern, z.\,B. auf Vakuummatratze: Trage mit
    vorbereiteter Matratze
    (\prettyref{topic:vakuummatratze-anwendung}) neben
    Patienten stellen um diesen einen möglichst geringen Weg
    tragen/bewegen zu müssen.
    \item Gurte öffnen und entfernen.
    \item Hälften trennen: beide Schnappverschlüsse öffnen
    \item Hälften parallel zum Patienten flach wegziehen. Der Patient soll
    hierbei nicht bewegt werden!
    \item Sollten die Schnappschlösser klemmen, kann man die
    Trage wie bei der Scherentechnik an nur einer Seite
    öffnen und in V-Form (gleichzeitig, je ein Helfer auf
    einer Seite) entfernen.
\end{enumerate}
% \end{multicols}
}{}{}{}{}
\end{ProcedureStepsNoHeading}

\subsubsection{\HeadingKeyword{Immobilisation} auf der Schaufeltrage}
\label{topic:schaufeltrage-immobilisation}

Auf manchen Schaufeltragen kann ein Patient direkt (ohne
Vakuummatratze) immobilisiert werden. Die Immobilisation auf der 
Schaufeltrage wird jedoch nicht in allen Rettungsdienstbereichen
empfohlen, da die aufnehmenden Spitäler i.\,d.\,R. nicht
darauf vorbereitet sind (Tauschmaterial, \ldots). \hfill
\Commentary[Schaufeltrage / Immobilisation]{Die
Immobilisation direkt auf der Schaufeltrage ist -- entgegen
vieler Gerüchte -- nicht nur möglich, sondern auch von
Herstellern, wie zum Beispiel Ferno{\texttrademark}
ausdrücklich freigegeben (z.\,B. Ferno
\E{ScoopEXL{\texttrademark}} Schaufeltrage). Die Produkte
sind auch entsprechnd konstruiert (Röntgendurchlässigkeit, \ldots).
Zubehör, wie z.\,B. die Ferno
\E{SpiderStraps{\texttrademark}} oder Ferno \E{Head
Immobilizer for 65EXL{\texttrademark}}, werden ausdrücklich
für Schaufeltragen angeboten.

Ein entsprechendes Instruktionsvideo kann unter
\url{http://www.ferno.com/scoopexl/scoopexl.htm} abgerufen
werden.  }

{\TakeNotesBoxLarge}


\begin{ProcedureStepsNoHeading}
\Text{Material und Personal}{
\begin{itemize}
    \item Schaufeltrage mit Zulassung zur Immobilisation
    von Patienten
    \item Passende Kopfstützen
    \item Passendes Gurtsystem zur Immobilisation (z.\,B.
    Ferno \E{Fastrap{\texttrademark} Quick Restraint
    System}, \FullPrettyRef{topic:spiderstraps-beschreibung})
    \item Patient auf Schaufeltrage
    (\FullPrettyRef{topic:schaufeltrage-aufschaufeln})
\end{itemize}
}{}{}{}{}




% \begin{multicols}{2}
\Text{Vorgehen}{
\begin{enumerate}
  \item Patient liegt, wie unter
  \FullPrettyRef{topic:schaufeltrage-aufschaufeln},
  beschrieben, auf der Schaufeltrage.
    \item Gurtsystem und Kopfstützen wie vom Hersteller
    beschrieben anbringen. 
    
    Je nach Schaufeltrage und Gurtsystem kann die Technik
    erheblich abweichen.
\end{enumerate}
% \end{multicols}

\begin{Danger}
  \item Bevor der Kopf fixiert wird, muss der Körper bereits
  verläßlich fixiert sein!
\end{Danger}
}{}{}{}{}
\end{ProcedureStepsNoHeading}

%%%%%%%%%%%%%%%%%%%%%%%%%%%%%%%%%%%%%%%%%%%%%%%%%%%%%%%%%%%
% TODO Sandwich-Technik: Neuerstellung
\TODO{GABS}{yyyy-mm-dd}{Sandwich-Technik: Neuerstellung von Grafik/Text 
                          notwendig}{}
%%%%%%%%%%%%%%%%%%%%%%%%%%%%%%%%%%%%%%%%%%%%%%%%%%%%%%%%%%%

% \clearpage % TEMP temp clearpage
\subsection{\HeadingKeyword{Vakuummatratze}}
\label{topic:vakuummatratze-anwendung}\index{Vakuummatratze!Handhabung}
\label{topic:vakuummatratze-beschreibung}




\Text{\TxBeschreibung}{
Die Vakuummatratze dient zur Ruhigstellung des gesamten
Körpers bei Verdacht auf mögliche Wirbelsäulenverletzungen,
Becken- oder Oberschenkelfrakturen (Oberschenkelhals und
Schaft).
Ziel ist es, die Matratze an die Körperform der Patienten
eng anzupassen und somit Bewegungen effektiv zu verhindern.
Vielerorts wird die Immobilisation mittels Vakuummatratze
als Mittel der Wahl empfohlen.
\hfill\Commentary[Vakuummatratze / Immobilisation]
{Die Immobilisation auf der Vakuummatratze wird derzeit als 
Methode der Wahl empfohlen, da die Technik
organisationsübergreifend weit verbreitet und das
Personal entsprechend geschult ist. Aufnehmende Spitäler
sind auf Patienten, welche auf einer Vakuummatratze
immobilisiert wurden, gut vorbereitet.}

{\TakeNotesBoxLarge}
}{}{}{}{}




\Picture
{Vakuummatratze}%1 Titel
{}%2 Beschreibung
{}%3 Label
{Ch. Pallinger}%4 Quelle
{MfG}%5 Lizenz
{./images/pallinger-christoph-aass/Vacuummatratze_32960-AASS-0176mm}%6 Datei
{}%7 Breitenmodifizierer
{}%8 Float
% vormals \Layout{37}

\begin{ProcedureStepsNoHeading}
\Text{Material und Personal}{
\begin{compactitem}
    \item Mind. 2 Helfer
    \item Trage
    \item Vakuummatratze
    \item Absauggerät
    \item Leintuch
    \item Evtl. Adapter, je nach Matratzentyp
\end{compactitem}
}{}{}{}{}


% \begin{multicols}{2}
\Text{Vorgehen}{
\begin{enumerate}
    \item Fahrtrage auf tiefste Stufe stellen (ganz unten).
    \item Vorbereiten: Vakuummatratze auf Fahrtrage legen. Die weichere
    Seite ist die Patientenseite, Ventil beim Kopf. Kugeln so verteilen,
    dass der gesamte Körper gut, sowie die betroffene Stelle besonders gut
    fixiert werden kann.
    \item Vorabsaugen mit Absauggerät bis die Matratze gut formbar aber nicht allzu
    fest ist.
    \item Vorformen, Ventil verschließen. Wenn vorhanden Leintuch auf
    die Matratze legen.
    \item Angegurteten Patienten mit Schaufeltrage auf die Vakuummatratze
    legen, Schaufeltrage entfernen (\prettyref{topic:schaufeltrage-handhabung}).
    \item \EE{Absaugen} mit Absauggerät, dabei zügig Matratze an den Patienten
    \EE{anpassen}. Besonderes Augenmerk gilt der Fixierung der verletzten
    Regionen sowie dem Kopf (\enquote{Ohren} formen: Ecken
    neben dem Kopf hoch formen). Einige Modelle bieten die
    Möglichkeit einer zusätzlichen Fixierung mittels an der
    Vakuummatratze befestigter Gurte.
    
    \begin{Danger}
      \item \Ca{Vorsicht!} Der Kopf und die HWS dürfen beim Anpassen nicht gestaucht werden!
    \end{Danger}
    \item Ist kein weiteres Absaugen mehr möglich und die
    Matratze gut angepasst (Matratze muss jetzt bretthart
    sein) Ventil schließen, Absauggerät abdrehen.
    \item Patient und Vakuummatratze an Fahrtrage angurten.
\end{enumerate}
% \end{multicols}
}{}{}{}{}
\end{ProcedureStepsNoHeading}

\subsection{\HeadingKeyword{Spineboard}}
\label{topic:spineboard-anwendung}\index{Spineboard!Anwendung|Z}
\label{topic:spineboard-beschreibung}\index{Spineboard!Beschreibung|Z}

\Text{\TxBeschreibung}{
Das Spineboard ist im Prinzip ein flaches Brett auf dem
der Patient ähnlich dem Schaufeltrage-Vakuummatratze-Duett
gerettet und fixiert werden kann. Es stellt somit in vielen
Situationen eine Alternative zur
Ganzkörper-/Wirbelsäulenimmobilisation dar.

Diese Technik wurde
ursprünglich vor allem im anglo-amerikanischen Raum
eingesetzt und erfährt auch hierzulande immer größere
Verbreitung.

}{}{}{}{}

\Customized{
\begin{description}
  \item[\CustAsboe] Zusätzlich zu den gebräuchlichen
  Hilfsmitteln zur Menschenrettung und Immobilisation kann
  ein
  Spineboard (Wirbelsäulenbrett) zur Rettung und
  Ganzkörperimmobilisierung traumatisierter Personen
  alternativ verwendet werden, sofern das entsprechende Material 
  vorhanden ist und das Personal darauf eingeschult wurde.
  \cite{AsboeLehrmeinungRd2011-10}
\end{description}
}
{\TakeNotesBoxLarge}




\PictureWide
{Spineboard}%1 Titel
{}%2 Beschreibung
{}%3 Label
{Ch. Pallinger}%4 Quelle
{MfG}%5 Lizenz
{./images/pallinger-christoph-aass/Spineboard_32956_crop-AASS-0112mm.jpg}%6 Datei
{}%7 Breitenmodifizierer
{}%8 Float
% vormals \Layout{37}

% \TakeNotesMedium

\subsection{\HeadingKeyword{Gurtensysteme} für Schaufeltrage und Spineboards}

\subsubsection{Ferno Fastrap{\texttrademark} Quick Restraint System}
\index{SpiderStraps!Beschreibung}
\label{topic:spiderstraps-beschreibung}
\index{Fastrap{\texttrademark}!Beschreibung}
\label{topic:fastrap-beschreibung}

\Text{\TxBeschreibung}
{%
% Anwendung: \prettyref{topic:schaufeltrage-anwendung}
Das Fastrap{\texttrademark}-System kann sowohl zusammen mit
einer Schaufeltrage, als auch mit einem Spineboard verwendet
werden. \Commentary[Ferno Fastrap{\texttrademark} /
Schaufeltrage]{\Speech{\enquote{The Fastrap system is compatible
with Ferno backboards and Scoop stretchers.}} Ferno International Online Shop,
\url{http://www.fernointernational.com}, 2011-05-13}
% 
% {\TakeNotesLarge}
}{

}
{}
{}
{}




\Picture
{Ferno Fastrap{\texttrademark} Quick Restraint System}%1 Titel
{}%2 Beschreibung
{}%3 Label
{Ch. Pallinger}%4 Quelle
{MfG}%5 Lizenz
{./images/pallinger-christoph-aass/Spiderstrap_32954_v2-AASS-0112mm.jpg}%6 Datei
{}%7 Breitenmodifizierer
{}%8 Float
% vormals \Layout{27}








% 
\subsection{\HeadingKeyword{Rettungskorsett}}
\label{topic:rettungskorsett-anwendung}\index{Rettungskorsett!Anwendung}\index{KED-System!anlegen}
\index{Rettungskorsett!Beschreibung}\label{topic:rettungskorsett-beschreibung}

\HeadingSub{Gängige Marken: Kendrick Extrication
  Device{\texttrademark} (KED{\texttrademark})}



\Text{\TxBeschreibung}
{%
Das Rettungskorsett wird primär für das Umlagern von
sitzenden Patienten mit dem Verdacht auf eine Verletzung der
Wirbelsäule verwendet. Am Kopfgriff des
Korsetts kann, falls nötig, ein Kranhaken eingehängt
werden. Das Rettungskorsett wird (ausser wenn es
ausschliesslich dazu verwendet wird, den Patienten 
zur Erleichterung der Bergung mit
\enquote*{Haltegriffen} auszustatten) gemeinsam mit einer
\EE{HWS-Schiene} verwendet.
}{

}
{}
{}
{}




\Picture
{Rettungskorsett}%1 Titel
{}%2 Beschreibung
{}%3 Label
{Wikipedia, \enquote{Onthost}}%4 Quelle
{PD}%5 Lizenz
{./images/geraete/KED.jpg}%6 Datei
{}%7 Breitenmodifizierer
{}%8 Float
% vormals \Layout{27}

\begin{ProcedureStepsNoHeading}
\Text{Material und Personal}{
\begin{compactitem}
    \item 2\,--\,3 Helfer
    \item Rettungskorsett
    \item HWS-Schiene
    \item Vorbereitete Vakuummatratze
\end{compactitem}
}{}
{}
{}
{}


% \begin{multicols}{2}
\Text{Vorgehen}{
\begin{enumerate}
  \item Patient aufsuchen und sofort HWS-Schiene anlegen.
  (\prettyref{topic:hws-schiene-anlegen})
  \item Ein Helfer stabilisiert von nun an am besten von
  hinten den  Kopf um seitliche Bewegungen des Kopfes zu
  vermeiden.
  \item Rettungskorsett \E{vorsichtig} in aufgefaltenem Zustand
  hinter dem Patienten platzieren. Die Oberkanten der
  Flügel sollen in der Achsel des Patienten zu liegen
  kommen.
  \item Die Beingurte vom Korsett lösen und baumeln
  lassen.
  \item Der mittlere und untere Bauchgurt wird geschlossen
  und halbfest angezogen. Das Korsett soll in der Höhe
  nicht mehr
  verrutschen können.
  \item Kopf mit der zum System gehörigen faltbaren
  Schaumstoffunterlage so unterpolstern, dass die HWS in
  Neutralposition steht.
  \item Der Kopf wird mit den Klettbändern fixiert: ein
  Band an der Stirn, das andere unter dem Kinn die HWS-Schiene
  entlang. Die Klettbänder werden am Rettungskorsett
  fixiert. Der Kopf soll nun seitlich nicht mehr geneigt
  werden können. Der Patient wird \E{unter keinen
  Umständen zu einem Versuch aufgefordert!}
  \item Die Beingurte werden geschlossen: von unten um
  den Oberschenkel herum nach oben; vorsichtig hin- und
  zurückziehend so eng es geht anlegen und festzurren.
  Vorsicht auf den Genitalbereich!
  \item Die halbfest geschlossenen Bauchgurte nun
  festzurren.
  \item Den letzten, obersten Bauchgurt schließen und
  festzurren. Dieser wird erst jetzt geschlossen um die
  Atemarbeit des Patienten nicht zu behindern.
  \item Der Patient ist nun entlang der gesamten
  Wirbelsäule geschient und kann mittels der 3 am Korsett
  angebrachten Griffe gehoben werden.
  \item Sanftes Umlagern auf die vorbereitete
  Vakuummatratze (\prettyref{topic:vakuummatratze-anwendung}).
  \item Beingurte leicht lockern.
  \item Den obersten Bauchgurt öffnen.
  \item Vakuummatratze anformen und sichern
  (\prettyref{topic:vakuummatratze-anwendung}).
  
\end{enumerate}
% \end{multicols}
}{}{}{}{}
\end{ProcedureStepsNoHeading}

\subsection{\HeadingKeyword{Vakuumschiene}}
\label{topic:vakuumschiene}
\label{topic:vakuumarmschiene}
\index{Vakuumarmschiene}

Analog zum Prinzip der Vakuummatratze gibt es auch
Vakuumschienen, um Extremitäten zu immobilisieren.
Gebräuchlich sind Vakuum-Bein- und Unterarmschienen, wobei
letztere kaum eingesetzt werden, da sie gegenüber eines
Dreiecktuchverbandes wenig Vorteile bieten.

\subsubsection{\HeadingKeyword{Vakuumbeinschiene}}
\label{topic:vacuumbeinschiene-anwendung}
\index{Vakuumbeinschiene!Anwendung}
\label{topic:vacuumbeinschiene-beschreibung}

\Text{\TxBeschreibung}
{%
Die Vakuumbeinschiene dient zum Ruhigstellen von
Verletzungen \EE{unterhalb des Knies}. Verletzungen des
Oberschenkels können mit ihr \E{nicht} ruhiggestellt werden,
da sie das Hüftgelenk nicht schient
(\Sign{→}\,Vakuummatratze)!
% 
% {\TakeNotesLarge}
}{

}
{}
{}
{}

\begin{ProcedureStepsNoHeading}
\Text{Material und Personal}{
\begin{compactitem}
    \item 2 Helfer: Einer seitlich, einer beim Fußende
    \item Vakuumbeinschiene
    \item Absauggerät
    \item Evtl. Adapter
\end{compactitem}
}{}{}{}{}



\PictureWide
{Vakuumbeinschiene mit Zubehör}%1 Titel
{}%2 Beschreibung
{}%3 Label
{Ch. Pallinger}%4 Quelle
{MfG}%5 Lizenz
{./images/pallinger-christoph-aass/Beinschiene_32977_v2-AASS-0176mm.jpg}%6 Datei
{}%7 Breitenmodifizierer
{}%8 Float
% vormals \Layout{37}


% \begin{multicols}{2}
\Text{Vorgehen}{
\begin{enumerate}
  \item Länge anpassen, gegebenenfalls oben umfalten.
  Klettverschlussbänder herrichten.
  \item Kugeln gut verteilen. Auch am Fußende müssen genug Kugeln
  sein, um dem Fuß sowohl unter der Sohle als auch am Fußrücken
  gut Halt geben zu können
  \item Schiene leicht vorabsaugen, damit Kugeln nicht mehr verrutschen
  können
  \item Schuh am gesunden Bein ausziehen: Finden der
  schonensten Methode.
  Ein Helfer hält das Bein auf Zug, der andere öffnet den
  Schuh (Bänder ausfädeln, zur Not aufschneiden)
  \item vorsichtiges Wiederholen zu zweit am verletztem Bein, Socken
  ausziehen (DMS)
  \item Stiefelgriff: gleich nach Ausziehen von Schuh/Socken
  nimmt der Helfer am Fußende durch das Loch in der Schiene
  die Ferse des verletzten Fußes, die andere Hand hält den
  Fußrücken (Stiefelgriff).
  \EE{Das Bein wird so auf dem Boden liegend auf Zug
  gehalten, der Helfer verbleibt so, bis die Schiene fest
  sitzt.}
  \item Der seitliche Helfer richtet die Schiene vorsichtig unter dem
  Patienten aus und fixiert die Schiene mit den
  Klettverschlussbändern.
  Um den Fuß werden Ohren gebogen und kreuzweise
  festgeklettet, alternativ ein Band über den Fußrücken und
  eines unter der Sohle. Der Fuß soll von allen Seiten so
  gut wie möglich umschlossen sein!
  \item Der seitliche Helfer saugt ab und zieht dabei die Klettbänder
  immer wieder fest nach. \EE{Der andere Helfer hält nach wie vor den
  Fuß im Stiefelgriff.}
  \item Fertig abgesaugt: Ventil schließen, Absauggerät abschalten. Der
  Stiefelgriff kann jetzt losgelassen werden.
  \item Nochmalige Kontrolle von Durchblutung, Motorik und
  Sensibilität (DMS) an der betroffenen Extremität
  (\FullPrettyRef{topic:dms}).
\end{enumerate}
% \end{multicols}
}{}{}{}{}
\end{ProcedureStepsNoHeading}

Die Schiene soll bretthart sein, das Bein leicht auf Zug
fixiert. Der Fuß soll weder seitlich noch zum Körper hin/vom
Körper weg beweglich sein.


\subsection{\HeadingKeyword{Aluminiumkernschiene}}
\label{topic:sam-splint-anwendung}
\label{topic:aluminiumkernschiene}
\label{topic:schienung-arm}
\index{Sam-Splint{\protect\texttrademark}!Anwendung}
\index{Aluminiumkernschiene!Anwendung}
\label{topic:sam-splint-beschreibung}
\index{Aluminiumkernschiene!Beschreibung}
\index{Sam-Splint{\protect\texttrademark}!Beschreibung}

\HeadingSub{Gängige Marken: SAM-Splint{\protect\texttrademark}}


\Text{\TxBeschreibung}
{%
Die Aluminiumkernschiene ist eine flexible, von Kunststoff
umgebene Metallschiene und ermöglicht das Ruhigstellen von
Extremitätenverletzungen mit abnormer Stellung. Sie lässt
sich leicht an den Arm anformen, muss jedoch noch mit einer
Mullbinde oder mit Dreiecktüchern fixiert werden. Die
\EE{Anpassung} erfolgt an der \EE{gesunden} Extremität.
}

\begin{PictureSeries}{}{Aluminiumkernschiene Marke Sam-Splint{\texttrademark}. \SourcePicture{Ch. Pallinger}}
\PictureSeriesRow{2}
{./images/pallinger-christoph-aass/Sam_Splint_33050_crop-AASS-0112mm}
{Ein SAM-Splint{\texttrademark} fabriksneu \ldots}
{./images/pallinger-christoph-aass/Sam_Splint_33053_crop-AASS-0112mm}
{{\ldots} und nach der Anpassung}
{}{}{}{}
\end{PictureSeries}

\begin{ProcedureStepsNoHeading}
\Text{Material und Personal}{
\begin{compactitem}
    \item 2 Helfer
    \item Sam-Splint
    \item Peha-Haft, Mullbinde + Leukoplast oder
    \VonBis{2}{3} Dreiecktuchkrawatten
    \item Ein weiteres Dreiecktuch als Armtragetuch
\end{compactitem}
}{}{}{}{}


\Text{Vorgehen}{
\begin{enumerate}
  \item Sam-Splint in der Hälfte falten, an gesundem Arm der Länge
  nach anpassen, Knick beim Ellbogen
  \item Anlegen am verletzten Arm, sanft anpassen. Fixierung mit
  Peha-Haft, Mullbinden oder Dreiecktuchkrawatten
  \item Fingerspitzen müssen sichtbar sein (DMS), gegebenenfalls Enden
  umfalten.
  \item Lagerung des geschienten Armes mittels Armtragetuch
\end{enumerate}
}{}{}{}{}
\end{ProcedureStepsNoHeading}

% 
% \needspace{0.5\textheight}
% ~

\section{\HeadingKeyword{Verbände} dienen dem Schutz und der Blutstillung}

\subsection{Druckverband}
\label{topic:druckverband}\index{Druckverband!anlegen}

\Text{\TxBeschreibung}{%
Der Druckverband dient der Stillung einer starken Blutung.
}{}{}{}{}

\begin{ProcedureStepsNoHeading}
\Text{Material und Personal}{%
\begin{compactitem}
    \item \VonBis{1}{2} Helfer
    \item Sterile Wundauflage je nach Wundgröße
    \item Druckkörper: z.\,B. Mullbinde
    \item Verbandsmaterial (Dreiecktuchkrawatte)
\end{compactitem}
}{}{}{}{}


% \begin{multicols}{2}
\Text{Vorgehen}{
\begin{enumerate}
    \item Sofort bei Erkennen einer \E{starken} Blutung
    durch direkten Druck auf die Wunde und Hochhalten die
    Blutung verlangsamen (Der Patient kann dies je nach Zustand
    unter Anleitung selbst tun).
    \item Sterile Wundauflage auf die Wunde legen.
    \item Druckkörper dem Wundverlauf folgend über die sterile
    Wundauflage legen, weiterhin drücken.
    \item Mit einem möglichst starren Verband\footnote{Druckverband: Unter Umständen ist die Verwendung von
    flexiblen Verbandsmaterial (z.\,B.
    PehaHaft{\texttrademark}) ausreichend bzw. je nach Wundlokalisation zu bevorzugen.} (z.\,B.
    Dreiecktuchkrawatte) wird der Druckkörper fest und
    möglichst vollständig umschlungen. 
    \item Das Verbandsmaterial wird festgezogen und mit einem festen Knoten und Gegenknoten bei
    ausreichendem Druck direkt über der Wunde fixiert.
    \item Sollte die Blutung zu einem späteren Zeitpunkt wieder
stärker werden oder der Druckverband zu locker angelegt
sein, wird ein \EE{weiterer} Druckkörper mit einer
weiteren Dreiecktuchkrawatte noch fester angelegt. \EE{Ein
schon angelegten Druckverband wird nicht wieder geöffnet!}
\end{enumerate}
% \end{multicols}
}{}{}{}{}
\end{ProcedureStepsNoHeading}

% \begin{PictureSeriesWide}{}{Bilderserie: \StepHeading{Druckverband} \SourcePicture{ASBÖ/gjh}}
% \PictureSeriesRowWide{3}
% {./images/ASBOE-Grassl-gjh-1/Baustelle_05-00441pt}
% {Die wichtigste Maßnahme bei einer blutenden Wunde ist das Hochlagern und Abdrücken unter Beachtung des Selbstschutzes (Schutzhandschuhe).}
% {./images/ASBOE-Grassl-gjh-1/Verband_02-00441pt}
% {Es folgt die Anlage eines Druckverbandes: Erst wird eine sterile Aundauflage aufgelegt, {\ldots}}
% {./images/ASBOE-Grassl-gjh-1/Verband_03-00441pt}
% {{\ldots} darüber ein Druckkörper.}
% {}
% {}
% \PictureSeriesRowWide{3}
% {./images/ASBOE-Grassl-gjh-1/Verband_05-00441pt}
% {Mit einem möglichst starren Verband (z.\,B. Dreiecktuchkrawatte) wird der Druckkörper umschlungen. U.\,U. ist auch flexibles Verbandsmaterial (z.\,B. PehaHaft{\texttrademark}) ausreichend.}
% {./images/ASBOE-Grassl-gjh-1/Verband_07-00441pt}
% {Anschließend wird der Knopf oberhalb der Wunde geknüpft.}
% {./images/ASBOE-Grassl-gjh-1/Baustelle_09-00441pt}
% {Danach wird die Verletzte Extremität sofern möglich hochgelagert.}
% {}
% {}
% \end{PictureSeriesWide}


% TODO Grassl Foto: Bilderserie Druckverband




\ChapterEnd